\documentclass[12pt]{article}
\usepackage[utf8]{inputenc}
\usepackage[T1]{fontenc}
\usepackage[margin=1in]{geometry}
\usepackage{setspace}
\usepackage{xcolor}
\usepackage[colorlinks=true, urlcolor=blue!60!black, linkcolor=blue!60!black]{hyperref}
\setlength{\parindent}{0em}
\setlength{\parskip}{0.45em}
\renewcommand{\baselinestretch}{1.0}
\begin{document}
\begin{center}
{\Large \textit{Clinton v. City of New York}, 524 U.S. 417 (1998)\par}
\vspace{0.3em}
{\normalsize Case Brief --- One-Page Edition\par}
\end{center}

\textbf{Memory Jogger.} The line-item veto case: canceling parts of signed laws violates the Presentment Clause (\href{https://supreme.justia.com/cases/federal/us/524/417/}{524 U.S. 417}).

\textbf{Facts.} The 1996 Line Item Veto Act let the President cancel spending items and limited tax benefits within five days of signing. Clinton canceled a New York Medicaid tax waiver and a tax deferral benefiting farmers' cooperatives.

\textbf{Procedural History.} New York City entities and an Idaho potato cooperative won in the district court; expedited appeal followed.

\textbf{Votes.} Affirmed, 6--3: Stevens for the Court; Kennedy concurred; Scalia and Breyer dissented in part.

\textbf{Holding.} The plaintiffs had standing, and cancellation is a unilateral partial repeal of enacted statutes that violates Article I, \S~7.

\textbf{Reasoning.} Under \textit{Chadha}, repeal requires bicameralism and presentment. A veto comes before enactment and covers the whole bill; cancellation comes after and strikes a part. Kennedy tied separation of powers to liberty; Scalia saw a valid delegation of power to ``decline to spend''; Breyer said the President was executing the statute.

\textbf{Examples.} \emph{Same-Side:} a ten-day earmark cancellation power. \emph{Opposite-Side:} fast-track rescission votes. \emph{Fence-Sitter:} spending directives under ``sums not exceeding'' appropriations.

\textbf{Critique.} Critics ask whether the line between declining to spend and erasing text is form over substance.

\textbf{Key Quotation.} ``Failure of political will does not justify unconstitutional remedies\ldots\ Liberty is always at stake when one or more of the branches seek to transgress the separation of powers\ldots\ Abdication of responsibility is not part of the constitutional design.'' --- Justice Kennedy, concurring, 524 U.S. at 449, 450, 452.

\end{document}
